\documentclass[11pt,a4paper]{article}
\usepackage[T1]{fontenc}
\usepackage{amsmath,amssymb,amsthm}
\usepackage[margin=27mm]{geometry}
\usepackage[hidelinks,pdftitle={Twist Algebra: Integral Coordinates for Projective States},pdfauthor={A. M. Thorn},pdflang={en}]{hyperref}
\newtheorem{proposition}{Proposition}
\newcommand{\Z}{\mathbb Z}
\newcommand{\Q}{\mathbb Q}
\newcommand{\C}{\mathbb C}
\newcommand{\PP}{\mathbb P}
\newcommand{\OO}{\mathcal O}
\newcommand{\ph}{\varphi}
\newcommand{\Nm}{\operatorname{N}}
\newcommand{\tr}{\operatorname{tr}}
\title{Twist Algebra:\\Integral Coordinates for Projective States}
\author{A. M. Thorn\\\small\href{https://orcid.org/0009-0008-5463-278X}{ORCID 0009-0008-5463-278X}}
\date{18 November 2026}
\begin{document}
\maketitle
\begin{abstract}
We describe projective states over the fifth cyclotomic field by an
integral matrix--scalar pair. Equality and overlap require only exact
field arithmetic, without a normalizing square root. For primitive
coordinates, the norm of the Hermitian size is a positive integral ray
height. We prove the sharp bound that this height is at least the square
of the number of nonzero coordinates, and classify its equality case.
In two coordinates this gives exactly twelve rays of height at most
four. The same bound explains the rigidity of integral invertible
similitudes. The construction specializes classical projective geometry
and Schmidt height; it does not derive physical measurement frequencies.
\end{abstract}

\section{Coordinates and primitive rays}\label{sec:coordinates}
Fix $j=e^{2\pi i/5}$, $K=\Q(j)$, $\OO=\Z[j]$ and
$F=\Q(\ph)$, where $\ph=(1+\sqrt5)/2=-j^2-j^3$.
Conjugation sends $j$ to $j^4$; $\dagger$ denotes conjugate transpose.
The ring of integers $\OO$ has discriminant $125$, real subring
$\Z[\ph]$ and roots of unity
$\mu_{10}=\{\pm j^r:0\leq r<5\}$ \cite{washington}.
Each coordinate has four integer coefficients in the basis
$(1,j,j^2,j^3)$, with $j^4=-1-j-j^2-j^3$.

States are rays $[v]\in\PP^{d-1}(K)$, $d\geq1$, in a fixed basis.
A vector in $\OO^d$ is \emph{primitive} if its coordinates generate
the unit ideal. Every ray has a primitive representative, unique up
to a unit. Indeed, Minkowski's ideal-class bound is
$(4/\pi)^2(4!/4^4)\sqrt{125}<2$; every class therefore contains
an integral ideal of norm one, so $\OO$ is a principal ideal domain.
Clear denominators, then divide by a generator of the coordinate ideal.
If $w=av$ and both are primitive, $\OO=a\OO$ forces
$a\in\OO^\times$.

\section{An integral pair for a ray}\label{sec:pairs}
For nonzero $v,w\in\OO^d$, set $h(v,w)=v^\dagger w$ and
$H(v)=h(v,v)$. The two embeddings $\sigma_1(j)=j$ and
$\sigma_2(j)=j^2$ represent the conjugate pairs of embeddings of $K$.
Their restrictions to $F$ are denoted $\tau_1,\tau_2$.
Then $H(v)\in\Z[\ph]$ and
$\tau_r(H(v))=\sum_i|\sigma_r(v_i)|^2>0$ for $r=1,2$.
Store the homogeneous pair
\begin{equation}\label{eq:pair}
 (R_v,s_v)=(vv^\dagger,H(v)),\qquad
 \rho_v=R_v/s_v\in M_d(K).
\end{equation}
This is the familiar rank-one projector representation of projective
geometry \cite{brody}, with numerator and denominator kept integral.

\begin{proposition}\label{prop:projector}
The matrix $\rho_v$ is Hermitian, idempotent and of trace one.
It determines the ray, and
\begin{equation}\label{eq:equality}
 [v]=[w]\quad\Longleftrightarrow\quad s_wR_v=s_vR_w.
\end{equation}
The projective overlap is
\begin{equation}\label{eq:overlap}
 P([v],[w])=\frac{h(v,w)\overline{h(v,w)}}{H(v)H(w)}
 =\tr(\rho_v\rho_w)\in F.
\end{equation}
It is unchanged by nonzero scalars in $K$, and
$0\leq\tau_r(P)\leq1$ for both embeddings.
Moreover $P=0$ precisely when $h(v,w)=0$, and $P=1$ precisely
when the rays coincide.
\end{proposition}
\begin{proof}
Use $R_v^2=H(v)R_v$ and $\tr R_v=H(v)$.
The image of $\rho_v$ is $Kv$. Scaling by $a\in K^\times$ multiplies
both entries by $a\bar a$; cross multiplication gives
\eqref{eq:equality}. Multiplying the two rank-one matrices gives
\eqref{eq:overlap}. Its numerator is fixed by conjugation.
Cauchy--Schwarz gives both bounds. Zero overlap means zero inner
product; equality at one in an embedding
means proportional complex vectors; choosing a nonzero coordinate
makes the proportionality factor an element of $K$.
\end{proof}

For example, $v=(1,1)$ and $w=(1,j)$ give
\begin{equation}\label{eq:example}
 P([v],[w])=\frac{1+\ph}{4},\qquad
 \tau_1(P)=\frac{3+\sqrt5}{8},\qquad
 \tau_2(P)=\frac{3-\sqrt5}{8}.
\end{equation}
One exact element has two real readings.

\section{A sharp height bound}\label{sec:height}
For a primitive representative write $H(v)=a+b\ph$ and define
\begin{equation}\label{eq:height}
 \Delta([v])=\Nm_{F/\Q}(H(v))=a^2+ab-b^2.
\end{equation}
This positive integer is unchanged by units: replacing $v$ by $uv$
multiplies it by $\Nm_{K/\Q}(u)=1$, a positive integer unit.
Primitivity matters: for example $(2,0)$ represents $[1:0]$ but
its unreduced norm is $16$, whereas the ray height is $1$.

This is the relative Schmidt height of the line $Kv$, not a new
height construction. In the product over all four embeddings, the
two Euclidean lengths of each conjugate pair multiply to
$\tau_r(H(v))$; the coordinate-ideal norm is one for primitive
$v$. Thus the definition in \cite{poels}, equation (2.1), gives
exactly $\Delta$. Its absolute Euclidean normalization is
$H_2([v])=\Delta([v])^{1/4}$.

\begin{proposition}\label{prop:height}
If a primitive $v$ has $k$ nonzero coordinates, then
\begin{equation}\label{eq:bound}
 \Delta([v])\geq k^2.
\end{equation}
Equality holds precisely when all its nonzero coordinates are
$u\varepsilon_i$, for one unit $u\in\OO^\times$ and
$\varepsilon_i\in\mu_{10}$.
\end{proposition}
\begin{proof}
For nonzero coordinates put $x_i=|\sigma_1(v_i)|$ and
$y_i=|\sigma_2(v_i)|$. The norm of every nonzero algebraic integer
in $K$ is a positive integer. Hence
\begin{equation}\label{eq:bound-proof}
 \Delta=(\sum_i x_i^2)(\sum_i y_i^2)
 \geq(\sum_i x_i y_i)^2
 =\left(\sum_i\sqrt{\Nm_{K/\Q}(v_i)}\right)^2\geq k^2.
\end{equation}
Equality forces each coordinate norm to be one, so each coordinate
is a unit, and Cauchy--Schwarz forces all $x_i/y_i$ to agree.
Since $x_i y_i=1$, the $x_i$ agree and the $y_i$ agree.
The ratio of any two nonzero coordinates is then an algebraic
integer with all conjugates of modulus one, hence a root of unity.
For completeness, this last assertion is Kronecker's lemma: the
monic polynomials of all positive powers have bounded integer
coefficients and bounded degree, so there are finitely many such
polynomials and roots. Two powers coincide, giving a root of unity.
Conversely, for coordinates $u\varepsilon_i$ one has
$H(v)=k u\bar u$ and $\Delta=k^2$.
\end{proof}

In two coordinates the complete lowest layer is therefore
\begin{equation}\label{eq:twelve}
 \{[v]\in\PP^1(K):\Delta([v])\leq4\}
 =\{[1:0],[0:1]\}\ \cup\ \{[1:\varepsilon]:\varepsilon\in\mu_{10}\}.
\end{equation}
Primitivity gives height one on the two axes; the equality case
gives the ten distinct ratios at height four. Heights two and three
do not occur.

\section{Which integral transformations can mix rays?}\label{sec:transformations}
For $A\in GL_d(K)$, clearing a common denominator gives an integral
representative of its projective action. The pair updates by
$R\mapsto ARA^\dagger$, $s\mapsto\tr(ARA^\dagger)$.
Before computing the new ray height, make the vector primitive.

\begin{proposition}\label{prop:rigidity}
If $A\in GL_d(\OO)$ and $A^\dagger A=cI$, then $A=uD\Pi$,
where $u\in\OO^\times$, $D$ is diagonal with entries in
$\mu_{10}$, and $\Pi$ is a permutation matrix.
For $c=1$, all nonzero entries lie in $\mu_{10}$.
\end{proposition}
\begin{proof}
Each column has Hermitian size $c$, so $c\in\Z[\ph]$ is totally
positive. Taking determinants and then norms gives
$\Nm_{F/\Q}(c)^d=\Nm_{K/\Q}(\det A)=1$.
Every column is primitive, since the corresponding row of $A^{-1}$
gives a linear combination of its entries equal to one.
Its height is therefore one; Proposition \ref{prop:height} leaves
exactly one nonzero entry, a unit. Invertibility places these entries
in different rows. They have the same Hermitian size $c$, so their
ratios have all conjugates of modulus one and lie in $\mu_{10}$.
If $c=1$, the entries themselves satisfy Kronecker's lemma.
\end{proof}

Thus an invertible integral standard similitude cannot mix coordinates.
An integral matrix invertible only over $K$ can:
\begin{equation}\label{eq:hadamard}
 A=\begin{pmatrix}1&1\\1&-1\end{pmatrix},\qquad
 A^\dagger A=A^2=2I,\qquad \det A=-2.
\end{equation}
It maps $[1:0]$ to $[1:1]$ and has inverse $A/2$ over $K$.
Projectively it is an involution; no division by $\sqrt2$ is needed.
The return sends $(1,1)$ to $(2,0)$, illustrating the required
primitive reduction before evaluating height.

\section{Scope}\label{sec:scope}
The scalar $J=1+j^2=j\ph^{-1}$ obeys $J\bar J=\ph^{-2}$,
yet $[Jv]=[v]$. Its action on integer coefficients differs from its
trivial action on $K$-rays. Height and support use the fixed standard
basis. The construction covers $K$-rays, not every complex ray;
its weights do not derive measurement frequencies, independent trials
or physical dynamics.

\paragraph{Computational checks.}
The accompanying verification scripts check pairs, norms, heights
and overlaps on exact finite examples.

\begin{thebibliography}{9}
\bibitem{washington} L. C. Washington, \emph{Introduction to
Cyclotomic Fields}, 2nd edition, Springer, 1997.
\href{https://doi.org/10.1007/978-1-4612-1934-7}{Publisher DOI}.
\bibitem{brody} D. C. Brody and L. P. Hughston, \emph{Geometric
Quantum Mechanics}, Journal of Geometry and Physics 38 (2001),
19--53. \href{https://arxiv.org/abs/quant-ph/9906086}{arXiv:quant-ph/9906086}.
\bibitem{poels} A. Poels, \emph{The complex case of Schmidt's
going-down Theorem}, Monatshefte für Mathematik 184 (2017),
649--666, equation (2.1).
\href{https://www.apoels-math-u.net/pdf/note_going_down.pdf}{Author manuscript};
\href{https://doi.org/10.1007/s00605-017-1084-4}{publisher DOI}.
\end{thebibliography}
\end{document}
